\section{Objectifs}
\begin{itemize}
\item Passer entre formes développée, canonique et factorisée.
\item Résoudre \texttt{ax²+bx+c=0}, factoriser lorsque c'est possible et étudier le signe.
\item Choisir une forme adaptée à une équation, une inéquation, une optimisation ou une étude de variations.
\end{itemize}

\section{Cours}
\subsection{Trois formes}
Pour \texttt{a≠0}, \texttt{f(x)=ax²+bx+c}. La forme développée est utile pour identifier les coefficients. La forme factorisée \texttt{a(x-x_1)(x-x_2)} rend les racines et le signe visibles. La forme canonique \texttt{a(x-α)²+β} met en évidence le sommet.

Le BO demande de savoir déterminer la forme canonique dans des cas simples par complétion du carré. \textbf{Le calcul effectif de la forme canonique dans le cas général n'est pas un attendu.}

\subsection{Discriminant et racines}
\texttt{Δ=b²-4ac}.
\begin{itemize}
\item \texttt{Δ<0} : aucune racine réelle ;
\item \texttt{Δ=0} : une racine double \texttt{-b/(2a)} ;
\item \texttt{Δ>0} : deux racines \texttt{(-b-√Δ)/(2a)} et \texttt{(-b+√Δ)/(2a)}.
\end{itemize}

Si deux racines \texttt{x_1,x_2} existent, leur somme vaut \texttt{-b/a} et leur produit \texttt{c/a}.

\subsection{Signe}
Lorsque deux racines distinctes existent, \texttt{f} est du signe de \texttt{a} à l'extérieur de l'intervalle entre les racines et du signe opposé à l'intérieur. Avec une racine double, le signe ne change pas.

\subsection{Stratégies de factorisation}
Avant d'utiliser systématiquement \texttt{Δ}, chercher une racine évidente, une identité remarquable, l'absence du terme en \texttt{x}, ou deux racines entières repérables par somme et produit.

\section{Démonstration à travailler}
La résolution générale de l'équation du second degré doit pouvoir être justifiée à partir d'une transformation algébrique menant au discriminant.

\section{Erreurs fréquentes}
\begin{itemize}
\item Calculer systématiquement une forme canonique générale alors que ce n'est pas l'objectif du programme.
\item Oublier le coefficient \texttt{a} dans la forme factorisée.
\item Déduire le signe sans tenir compte du signe de \texttt{a}.
\end{itemize}

\section{Synthèse}
Le second degré se traite en choisissant la forme qui rend la question la plus simple : développée pour les coefficients, factorisée pour racines et signe, canonique pour le sommet et l'optimisation.
